\documentclass[11pt,a4paper]{article}

% ============================================================
% PACKAGES
% ============================================================
\usepackage[utf8]{inputenc}
\usepackage[T1]{fontenc}
\usepackage[french]{babel}
\usepackage{amsmath, amssymb, amsfonts}
\usepackage[margin=1.3cm, top=1.3cm, bottom=1.1cm]{geometry}
\usepackage{enumitem}
\usepackage{fancyhdr}
\usepackage{tikz}
\usetikzlibrary{arrows.meta}
\usepackage{multicol}
\usepackage{tcolorbox}
\tcbuselibrary{skins}

% ============================================================
% EN-TÊTE / PIED DE PAGE
% ============================================================
\setlength{\headheight}{14pt}
\pagestyle{fancy}
\fancyhf{}
\fancyhead[L]{\small\itshape Équation de droites}
\fancyhead[R]{\small\itshape Première STI2D}
\fancyfoot[C]{\small\thepage}
\renewcommand{\headrulewidth}{0.4pt}

% ============================================================
% LISTES COMPACTES
% ============================================================
\setlist[itemize]{leftmargin=12pt, itemsep=0pt, topsep=1pt, parsep=0pt}
\setlist[enumerate]{leftmargin=16pt, itemsep=2pt, topsep=1pt, parsep=0pt}

\setlength{\parindent}{0pt}

% ============================================================
% DOCUMENT
% ============================================================
\begin{document}

% -------- TITRE --------


\vspace{-4pt}
\begin{tcolorbox}[colback=gray!5, colframe=gray!40, boxrule=0.4pt,
                  arc=2pt, left=6pt, right=6pt, top=3pt, bottom=3pt]
\small
\begin{center}
{\large\bfseries Rappels de cours}\\[2pt]
{\small\itshape Équations de droites — Première STI2D}
\end{center}
\end{tcolorbox}
\vspace{4pt}

% ============================================================
% ================  ANNEXE — RAPPELS DE COURS  ===============
% ============================================================



\vspace{-4pt}
\hrule
\vspace{4pt}

\small
\begin{multicols}{2}

% ---- 1 ----
\noindent\textbf{1. Repère et coordonnées}

Dans un repère du plan, tout point $M$ est repéré par un unique couple
$(x_M\,;\,y_M)$ : $x_M$ est son \emph{abscisse}, $y_M$ son \emph{ordonnée}.

Un repère est \emph{orthonormé} (ou orthonormal) lorsque ses deux axes sont
perpendiculaires et que l'unité de longueur est la même sur les deux axes.

\medskip
% ---- 2 ----
\noindent\textbf{2. Équation réduite d'une droite}

Toute droite \emph{non verticale} admet une unique équation réduite de la forme
\[
y = m\,x + p,
\]
où $m$ est le \emph{coefficient directeur} et $p$ l'\emph{ordonnée à l'origine}.

Une droite \emph{verticale} a pour équation $x = k$ ($k \in \mathbb{R}$).

\medskip
% ---- 3 ----
\noindent\textbf{3. Coefficient directeur}

Le coefficient directeur mesure l'\emph{inclinaison} de la droite.

\emph{À partir de deux points :} si $A(x_A\,;\,y_A)$ et $B(x_B\,;\,y_B)$ sont
deux points distincts d'une droite non verticale ($x_A \neq x_B$), alors
\[
m = \frac{y_B - y_A}{x_B - x_A}.
\]

\emph{Interprétation :} en partant d'un point de la droite, lorsque l'on se
déplace de $1$ unité vers la droite, l'ordonnée varie de $m$ unités.

\medskip
% ---- 4 ----
\noindent\textbf{4. Ordonnée à l'origine}

Dans l'équation $y = m\,x + p$, le réel $p$ est l'ordonnée du point
d'intersection de la droite avec l'axe des ordonnées : ce point a pour
coordonnées $(0\,;\,p)$.

\medskip
% ---- 5 ----
\noindent\textbf{5. Équation cartésienne d'une droite}

Toute droite du plan admet une équation de la forme
\[
a\,x + b\,y + c = 0,
\]
où $a$, $b$, $c$ sont des réels, et $(a\,;\,b) \neq (0\,;\,0)$.
Si $b \neq 0$, on peut la mettre sous forme réduite
$y = -\dfrac{a}{b}\,x - \dfrac{c}{b}$.

\medskip
% ---- 6 ----
\noindent\textbf{6. Droites parallèles}

Deux droites \emph{non verticales} sont parallèles si et seulement si elles
ont le même coefficient directeur :
\[
d \parallel d' \iff m = m'.
\]

\medskip
% ---- 7 ----
\noindent\textbf{7. Droites perpendiculaires}

Deux droites \emph{non verticales} sont perpendiculaires si et seulement si
le produit de leurs coefficients directeurs est égal à $-1$ :
\[
d \perp d' \iff m \times m' = -1.
\]

\medskip
% ---- 8 ----
\noindent\textbf{8. Méthode : équation d'une droite passant par deux points}

\begin{enumerate}[leftmargin=14pt, itemsep=1pt]
\item Calculer le coefficient directeur $m = \dfrac{y_B - y_A}{x_B - x_A}$.
\item Écrire $y = m\,x + p$ et remplacer $x$ et $y$ par les coordonnées de l'un
des deux points pour déterminer $p$.
\item Conclure en donnant l'équation réduite $y = m\,x + p$.
\end{enumerate}

\medskip
% ---- 9 ----
\noindent\textbf{9. Méthode : équation d'une droite passant par un point et de
coefficient directeur connu}

\begin{enumerate}[leftmargin=14pt, itemsep=1pt]
\item Partir de la forme $y = m\,x + p$ avec le $m$ donné.
\item Remplacer $x$ et $y$ par les coordonnées du point $A$ pour déterminer $p$.
\item Conclure en donnant l'équation réduite.
\end{enumerate}

\medskip
% ---- 10 ----
\noindent\textbf{10. Méthode : équation d'une droite passant par un point et
d'ordonnée à l'origine connue}

\begin{enumerate}[leftmargin=14pt, itemsep=1pt]
\item Partir de la forme $y = m\,x + p$ avec le $p$ donné.
\item Remplacer $x$ et $y$ par les coordonnées du point $A$ pour déterminer $m$.
\item Conclure en donnant l'équation réduite.
\end{enumerate}

\medskip
% ---- 11 ----
\noindent\textbf{11. Méthode : droite parallèle à une droite donnée passant par un point}

\begin{enumerate}[leftmargin=14pt, itemsep=1pt]
\item Lire le coefficient directeur $m$ de la droite donnée.
\item Utiliser la propriété $d' \parallel d \iff m' = m$.
\item Déterminer $p'$ en utilisant les coordonnées du point imposé.
\end{enumerate}

\medskip
% ---- 12 ----
\noindent\textbf{12. Sens de variation et signe du coefficient directeur}

\begin{itemize}[leftmargin=12pt]
\item Si $m > 0$, la droite « monte » lorsque $x$ augmente.
\item Si $m < 0$, la droite « descend » lorsque $x$ augmente.
\item Si $m = 0$, la droite est horizontale (parallèle à l'axe des abscisses).
\end{itemize}

\medskip
% ---- 13 ----
\noindent\textbf{13. Vocabulaire de mise en situation}

Dans un problème concret, le coefficient directeur correspond souvent à une
\emph{pente}, un \emph{taux d'accroissement} ou une \emph{vitesse} :
\[
\text{pente} = \frac{\text{dénivelé}}{\text{distance horizontale}} = m.
\]

\end{multicols}


\end{document}